\documentclass[12pt,a4paper]{article}
\usepackage[margin=2.5cm]{geometry}
\usepackage{booktabs}
\usepackage{array}
\usepackage{amsmath}
\usepackage{amssymb}
\usepackage{listings}
\usepackage{xcolor}
\usepackage{titlesec}
\usepackage{enumitem}
\usepackage{tabularx}
\usepackage{multirow}
\usepackage{fancyhdr}
\usepackage{hyperref}

\lstset{
  basicstyle=\ttfamily\small,
  breaklines=true,
  frame=single,
  backgroundcolor=\color{gray!10},
  keywordstyle=\color{blue}\bfseries,
  language=SQL
}

\pagestyle{fancy}
\fancyhf{}
\lhead{02170 Database Systems}
\rhead{Written Examination 2023}
\cfoot{\thepage}

\titleformat{\section}{\large\bfseries}{}{0em}{}[\titlerule]
\titleformat{\subsection}{\normalsize\bfseries}{}{0em}{}

\begin{document}

%% ============================================================
%%  COVER PAGE
%% ============================================================
\begin{center}
  \vspace*{1cm}
  {\large \textbf{Technical University of Denmark}}\\[0.3em]
  Department of Applied Mathematics and Computer Science\\[1.5em]
  {\LARGE \textbf{02170 Database Systems}}\\[0.5em]
  {\Large Written Examination 2023}\\[1em]
  \hrule
  \vspace{0.5em}
  \begin{tabular}{ll}
    \textbf{Exam date:}     & May 2023 \\
    \textbf{Duration:}      & 3 hours \\
    \textbf{Aids:}          & All written materials permitted. No internet. \\
  \end{tabular}
  \vspace{0.5em}
  \hrule
\end{center}

\vspace{1em}
\noindent\textbf{Note:} The exam consists of two parts. Part 1 is SQL (written queries). Part 2 is multiple choice theory. The parts can be solved in any order.

\newpage

%% ============================================================
%%  SCENARIO
%% ============================================================
\section*{Scenario: A Database for an Airline}

A database named \texttt{Airline} is considered. The database contains information about flights offered by an airline company named \textit{Noob Airlines}, passengers on the flights, and customers. Below you can find the relation schemas for the five relations in the database, and a relation instance of each of them.

\vspace{0.5em}
\begin{tabular}{ll}
\textbf{Airport}(\underline{airportID}, city) &
\textbf{Person}(\underline{personID}, name, speed) \\
\end{tabular}

\vspace{0.3em}
\textbf{Plane}(\underline{planeID}, capacity)

\vspace{0.3em}
\textbf{Flight}(\underline{flightID}, dayTime, duration, fromAirport, toAirport, planeID, departureDate, latestDepTime)

\vspace{0.3em}
\textbf{FlightPassengers}(\underline{flightID}, \underline{customerID}, price)

\vspace{1em}
\noindent\textbf{Airport:} Each airport has a unique \textit{airportID} and the city in which the airport is situated.

\noindent\textbf{Plane:} Each plane has a unique \textit{planeID} and a capacity stating the number of passengers that can be seated in the plane.

\noindent\textbf{Flight:} Each flight has a unique \textit{flightID}, a departure date (\textit{dayTime}) and a departure time (\textit{latestDepTime}) (in the CET time zone), and a duration. It also has a departure airport (\textit{fromAirport}) and a destination airport (\textit{toAirport}). A specific plane (\textit{planeID}) can be allocated to the flight. Initially no plane has been allocated, therefore \textit{planeID} can be NULL.

\noindent\textbf{Customer:} Before a customer can be booked on a flight, the customer must be registered with a unique \textit{customerID} and a name.

\noindent\textbf{FlightPassengers:} In this table it is registered which customers are booked as passengers on which flights, and which individual price each passenger of a flight has paid for that flight.

\vspace{1em}

% Airport table
\begin{minipage}[t]{0.3\textwidth}
\textbf{Airport}\\[0.3em]
\begin{tabular}{|l|l|}
\hline
\textbf{airportID} & \textbf{city} \\\hline
CPH & Copenhagen \\\hline
MXP & Milano \\\hline
ARN & Stockholm \\\hline
OSL & Oslo \\\hline
LHR & London \\\hline
\end{tabular}
\end{minipage}
\hfill
% Plane table
\begin{minipage}[t]{0.3\textwidth}
\textbf{Plane}\\[0.3em]
\begin{tabular}{|l|l|}
\hline
\textbf{planeID} & \textbf{capacity} \\\hline
P1 & 180 \\\hline
P2 & 220 \\\hline
P3 & 150 \\\hline
\end{tabular}
\end{minipage}
\hfill
% Customer table
\begin{minipage}[t]{0.3\textwidth}
\textbf{Customer}\\[0.3em]
\begin{tabular}{|l|l|}
\hline
\textbf{customerID} & \textbf{name} \\\hline
C1 & Alice \\\hline
C2 & Bob \\\hline
C3 & Carol \\\hline
C4 & Dave \\\hline
\end{tabular}
\end{minipage}

\vspace{1em}

% Flight table
\textbf{Flight}\\[0.3em]
\begin{tabular}{|l|l|l|l|l|l|}
\hline
\textbf{flightID} & \textbf{dayTime} & \textbf{duration} & \textbf{fromAirport} & \textbf{toAirport} & \textbf{planeID} \\\hline
SK101 & 2023-05-01 08:00 & 120 & CPH & MXP & P1 \\\hline
SK102 & 2023-05-01 10:00 & 90  & MXP & CPH & P1 \\\hline
SK103 & 2023-05-02 09:00 & 110 & CPH & ARN & P2 \\\hline
SK104 & 2023-05-02 14:00 & 80  & ARN & CPH & NULL \\\hline
SK105 & 2023-05-03 07:00 & 130 & CPH & LHR & P3 \\\hline
SK106 & 2023-05-03 15:00 & 75  & OSL & CPH & P2 \\\hline
\end{tabular}

\vspace{1em}

% FlightPassengers table
\textbf{FlightPassengers}\\[0.3em]
\begin{tabular}{|l|l|l|}
\hline
\textbf{flightID} & \textbf{customerID} & \textbf{price} \\\hline
SK101 & C1 & 1200 \\\hline
SK101 & C2 & 1500 \\\hline
SK101 & C3 & 900  \\\hline
SK102 & C1 & 1100 \\\hline
SK103 & C2 & 800  \\\hline
SK103 & C4 & 1300 \\\hline
SK105 & C3 & 1250 \\\hline
SK105 & C4 & 1250 \\\hline
\end{tabular}

\newpage

%% ============================================================
%%  PART 1 — SQL
%% ============================================================
\section*{Part 1 — SQL Queries}

\noindent\textit{Write at most one SQL query per question. Place your answer in the answer file.}

\vspace{1em}

\subsection*{Question 1}
Write an SQL query that returns a table containing the \texttt{flightID} of those flights that are flying to \textbf{Milano} city.

\vspace{2cm}

\subsection*{Question 2}
Write an SQL query that returns a table containing the \texttt{flightID} of flights for which the highest price (paid by a passenger) is \textbf{less than 1300}.

\vspace{2cm}

\subsection*{Question 3}
Write an SQL query that returns a table containing the \texttt{customerID} and the \textbf{number of free seats} for each flight which goes \textbf{from the CPH airport} and which has got a plane allocated.

\textit{Hint: It is assumed that each passenger booked on a flight will be assigned exactly one seat on that flight. Furthermore, the airline sometimes overbooks flights; in such a case the number of free seats will be negative.}

\vspace{2cm}

\subsection*{Question 4}
Define an SQL trigger named \texttt{Flight\_Before\_Insert}, which automatically raises a signal when a new row having the \textbf{same \texttt{fromAirport} and \texttt{toAirport}} is attempted to be inserted in the \texttt{Flight} table.

\vspace{2cm}

\subsection*{Question 5}
Show an example of an SQL statement which \textbf{causes this trigger to raise a signal}.

\vspace{2cm}

\newpage

%% ============================================================
%%  PART 2 — MULTIPLE CHOICE
%% ============================================================
\section*{Part 2 — Multiple Choice Theory}

\noindent\textit{Select the single best answer for each question unless stated otherwise.}

\vspace{1em}

%% --- Q1: ER Modelling ---
\subsection*{Question 1: ER Modelling}

Consider the following relation schemas:

\begin{itemize}
  \item \textbf{Airport}(\underline{airportID}, city)
  \item \textbf{Flight}(\underline{flightID}, fromAirport, toAirport, planeID)
\end{itemize}

Which of the following ER diagrams correctly depicts the relationship between \texttt{Airport} and \texttt{Flight}?

\vspace{0.5em}
\noindent(Multiple ER diagram options — see original exam)

\vspace{1em}

%% --- Q2: Design ---
\subsection*{Question 2: ER Design}

The database is extended to also track which customers have a \textit{FrequentFlyerCard}. Consider the partial ER model below.

\vspace{0.5em}
\noindent\textbf{Question 2.1:} Which of the following statements is true / false?
\begin{enumerate}[label=\alph*)]
  \item Every customer has a FrequentFlyerCard.
  \item There are FrequentFlyerCards that are not assigned to any customer.
  \item A customer can have multiple FrequentFlyerCards.
  \item Every FrequentFlyerCard belongs to exactly one customer.
  \item There are customers that have been assigned to multiple cards.
\end{enumerate}

\vspace{0.5em}
\noindent\textbf{Question 2.2: Transformation ER to Relational Model}\\
Which of the following options is the correct transformation to the relational model?
\begin{enumerate}[label=\alph*)]
  \item \textit{[option a — see original exam]}
  \item \textit{[option b — see original exam]}
  \item \textit{[option c — see original exam]}
  \item \textit{[option d — see original exam]}
\end{enumerate}

\vspace{1em}

%% --- Q3: Normalization ---
\subsection*{Question 3: Normalization}

Consider a relation with schema $R(A, B, C, D, E)$ where the primary key is not shown.

The following functional dependencies are assumed:
$$FD = \{A \rightarrow BCD,\ CD \rightarrow E\}$$

\textbf{Question 3.1: Primary Key}\\
What is/are the candidate key(s) of $R$?

\vspace{1cm}

\textbf{Question 3.2: Normal Form}\\
What is the highest normal form of $R$?

\vspace{1cm}

\textbf{Question 3.3: Normalization}\\
Normalize $R$ into third normal form (3NF), if it is not already in 3NF.

\vspace{1cm}

%% --- Q4: Integrity Constraints ---
\subsection*{Question 4: Integrity Constraints}

\textbf{Question 4.1: Data Integrity}\\
Name two options to enforce that a user can never insert a new row into \texttt{Flight} where \texttt{fromAirport = toAirport}.

\vspace{1cm}

\textbf{Question 4.2: Foreign Key Actions}\\
Which referential action would you pick to ensure that a flight cannot be deleted from the database while it still has passengers booked on it (i.e., there is a row in \texttt{FlightPassengers} that references it)?

\vspace{1cm}

%% --- Q5: Indexing and Hashing ---
\subsection*{Question 5: Indexing and Hashing}

Assume that there is exactly one index on the \texttt{Flight} table (no automatic indexes were created), created via:

\begin{lstlisting}
CREATE INDEX flight_idx ON Flight(fromAirport) USING HASH;
\end{lstlisting}

Which of the following SQL queries is sped up by this index?
\begin{enumerate}[label=\alph*)]
  \item \texttt{SELECT * FROM Flight WHERE fromAirport <= 'CPH';}
  \item \texttt{SELECT * FROM Flight WHERE fromAirport = 'CPH' AND toAirport = 'MXP';}
  \item \texttt{SELECT * FROM Flight WHERE fromAirport > 'CPH' AND toAirport = 'MXP';}
  \item \texttt{SELECT * FROM Flight WHERE toAirport = 'MXP';}
  \item \texttt{SELECT * FROM Flight WHERE flightID = 'SK101';}
\end{enumerate}

\vspace{1em}

%% --- Q6: Formal Query Languages ---
\subsection*{Question 6: Formal Query Languages}

Which of the following options in domain calculus is the equivalent of the relational algebra statement below?

$$\Pi_{flightID}\left(\sigma_{city='Milano'}\left(Flight \bowtie Airport\right)\right)$$

\begin{enumerate}[label=\alph*)]
  \item $\{flightID \mid \langle flightID, fromAirport, toAirport \rangle \in Flight\ \wedge\ \langle toAirport, \text{'Milano'} \rangle \in Airport\}$
  \item $\{flightID \mid \exists fromAirport, toAirport\ (\langle flightID, fromAirport, toAirport \rangle \in Flight\ \wedge\ \langle toAirport, \text{'Milano'} \rangle \in Airport)\}$
  \item $\{flightID \mid \exists fromAirport, toAirport\ (\langle flightID, fromAirport, toAirport \rangle \in Flight)\ \wedge\ \exists toAirport\ (\langle toAirport, \text{'Milano'} \rangle \in Airport)\}$
  \item $\{flightID \mid \exists fromAirport, toAirport, city\ (\langle flightID, fromAirport, toAirport \rangle \in Flight\ \wedge\ \langle toAirport, city \rangle \in Airport)\}$
\end{enumerate}

\vspace{2em}
\begin{center}
\textit{--- End of Exam ---}
\end{center}

\end{document}
